% Source-aligned English for printed p.167.
\SourcePage{167}

\LayerTarget{tohoku:construction:cohomology-support-comparison:p2}
\noindent that in degree 0 reduce to the one considered above); alternatively,
starting from a sheaf $F$ on $X$,

\LayerTarget{tohoku:eq:canonical-support-comparison-FZ}
\[
H^p_{\Phi_Z}(Z,F\vert Z)\to H^p_\Phi(X,F_Z).
\]

\LayerTarget{tohoku:theorem:3.5.1}
\indent\textsc{Theorem 3.5.1.} \textit{The preceding homomorphisms are
isomorphisms in each of the following two cases:}
\begin{enumerate}
\item \textit{$Z$ is closed.}
\item \textit{$\Phi$ is paracompactifying and $Z$ is open.}
\end{enumerate}

\LayerTarget{tohoku:proof:3.5.1}
\indent\textsc{Proof.} In each case it is enough to verify that the functors
$H^p_\Phi(X,G^X)$ on $\mathbf C^Z$ are effaceable for $p\geqq1$. If $Z$ is closed,
this follows from the fact that $G^X$ is injective whenever $G$ is injective
(Prop. 3.1.4). If $Z$ is open and $\Phi$ is paracompactifying, then so is
$\Phi_Z$; hence every $G\in \mathbf C^Z$ embeds in a $\Phi_Z$-soft sheaf (cf. 3.3).
It remains only to note that if $G$ is $\Phi_Z$-soft, then $G^X$ is
$\Phi$-soft, which is immediate, and to apply the corollary to Proposition
3.3.2.

Return now to the exact sequence
$0\to F_U\to F\to F_Y\to0$
associated with a closed subspace $Y$ and its open complement $U$. It gives
rise to a cohomology exact sequence which, by the preceding theorem, may be
written

\LayerTarget{tohoku:eq:closed-open-support-long-exact-sequence}
\[
\mathinner{\cdotp\cdotp\cdotp\cdotp} H^p_\Phi(X,F_U)\to H^p_\Phi(X,F)
\to H^p_{\Phi_Y}(Y,F)\to
H^{p+1}_\Phi(X,F_U)\mathinner{\cdotp\cdotp\cdotp\cdotp}
\]
(where, for simplicity, $F$ has been written in place of $F\vert Y$ in the
third term).
\LayerTarget{tohoku:eq:open-support-identification}
If $\Phi$ is paracompactifying, terms of the form $H^p_\Phi(X,F_U)$ may
further be replaced by $H^p_{\Phi_U}(U,F)$, giving the familiar exact
sequence of [4, XVII].

\LayerTarget{tohoku:section:3.6}
\indent\textbf{3.6. On the cohomological dimension of certain spaces}%
\LayerTarget{tohoku:footnote:8}\textsuperscript{8)}.

\LayerTarget{tohoku:proposition:3.6.1}
\indent\textsc{Proposition 3.6.1.} \textit{Let $X$ be a topological
space and $(T^p)$ a covariant cohomological functor, defined on the category
$\mathbf C^X$ of sheaves of abelian groups on $X$ and taking values in a category
$\mathbf C'$. Suppose that $\mathbf C'$ satisfies condition \textup{AB 4)} \textup{(}cf. \textup{1.5)}, which implies
that inductive limits may be formed in $\mathbf C'$ \textup{(}Prop. \textup{1.8)}, and suppose that the
$T^p$ commute with inductive limits. If $F\in \mathbf C^X$, then $T^p(F)$ belongs to
every thick subcategory \textup{(}cf. \textup{1.11)} $\mathbf C''$ of $\mathbf C'$ that is stable under infinite
direct sums and contains every object of the form $T^i(\mathbf Z_U)$, where $U$ is an
arbitrary open subset of $X$ and $i$ is equal to $p,p+1$, or $p+2$.}

\LayerTarget{tohoku:proof:3.6.1:p1}
(The meaning of $\mathbf Z_U$ is the same as in the preceding section.)
\LayerTarget{tohoku:construction:generator-family}
Consider a family $(f_i)_{i\in I}$ of sections of $F$ on open subsets $U_i$.
Each $f_i$ defines a homomorphism from $\mathbf Z_{U_i}$ to $F$, and hence $(f_i)$
defines a homomorphism from the
\LayerTarget{tohoku:eq:generator-epimorphism}
direct sum $\displaystyle\bigoplus_i \mathbf Z_{U_i}$ to $F$.
\LayerTarget{tohoku:def:system-of-generators}
We call $(f_i)$ a \textit{system of generators} of $F$ if the preceding
homomorphism is an epimorphism. It is clear that, for any given $F$, a family
of generators always exists. Consequently, $F$ is the inductive limit of an
increasing filtered family

\TohokuContinuation{168}
